\section*{INTRODUCTION}

Animals swim, walk, and breathe in steady rhythms.
These rhythmic behaviors rely on central pattern generators (CPGs), neural circuits that can generate coordinated motor output in the absence of patterned sensory input.
Classic fictive preparations established that the temporal structure of motor patterns can arise centrally, motivating the study of CPGs as autonomous yet context-sensitive networks \cite{Wilson1961, MarderCalabrese1996, Kiehn2006}.
Early work on invertebrate CPGs recognized two broad architectures: pacemaker-driven oscillators, where intrinsically bursting neurons provide rhythmic drive, and network oscillators, where rhythm emerges from synaptic interactions among non-bursting neurons \cite{Getting1989}.
The Tritonia swim CPG exemplified the latter---none of its neurons exhibited intrinsic bursting, yet the network generated robust rhythm through ``delayed excitation'' and other slow synaptic mechanisms \cite{Getting1983delayed, Getting1983network}.
As research characterized diverse ionic currents underlying cellular bursting \cite{Rinzel87b, RL87}, attention shifted toward pacemaker-driven systems and the modulation of intrinsic properties.
The question of how network oscillators achieve \textit{smooth, graded frequency control} has received less systematic attention.

The swim CPG of the nudibranch \textit{Dendronotus iris} presents a compact instance of this problem: four neurons, none capable of intrinsic bursting, generate robust left--right alternation whose frequency varies smoothly under a single command signal.
Like its Tritonia relatives, the \textit{Dendronotus} circuit relies on synaptic rather than cellular mechanisms for rhythm generation.
We propose that such circuits exploit what we term \textit{network hysteresis}: slow synaptic dynamics assume the role of the slow intrinsic variable found in cellular bursters, accumulating across network partners to create history-dependent transitions between collective activity states.
This terminology builds on Getting's (1983) characterization of ``delayed excitation,'' providing a unifying concept that connects slow synaptic accumulation to the hysteresis framework used to analyze cellular bistability.

Intrinsically bursting neurons generate rhythm through cellular hysteresis: active and silent states coexist across a range of a slow intrinsic variable such as intracellular calcium, and the variable's drift between these states drives phase transitions within a single cell \cite{Rinzel87b, RL87}.
This mechanism imposes rhythm but couples frequency tightly to burst shape.
Half-center oscillators enforce alternation through reciprocal inhibition \cite{brown}, with escape, release, or post-inhibitory rebound governing phase transitions \cite{Wang-Rinzel-92, SKM94}.
Skinner, Kopell, and Marder (1994) showed that when phase transitions depend on intrinsic cellular mechanisms, network frequency is insensitive to synaptic parameters; when transitions depend on synaptic mechanisms, frequency becomes tunable through synaptic properties.
This distinction suggests that network oscillators lacking cellular pacemakers may be inherently more amenable to frequency control---a possibility we explore quantitatively.

Network oscillators can be understood as assemblies of canonical subnetwork building blocks \cite{Getting1989}.
Excitatory-inhibitory (E/I) coupling represents one such building block capable of generating oscillations through network interaction alone \cite{WilsonCowan1972}.
Wilson and Cowan (1972) showed that E/I networks can exhibit oscillations and hysteresis, but their models operate on cortical timescales with fast synapses tuned for gamma rhythms.
We propose that network hysteresis, arising from slow synaptic dynamics within E/I subnetworks, provides the mechanism by which network oscillators achieve robust, graded frequency control at CPG timescales.

CPG circuits exhibit degeneracy: similar output despite varied synaptic and cellular properties \cite{MarderCalabrese1996, Marder2011}.
Neuromodulators reconfigure circuits without changing wiring \cite{NusbaumBeenhakker2002, JingWeiss2005}, but adaptability also demands that frequency be tunable.

Slow synaptic dynamics provide a plausible substrate for such control.
Repeated presynaptic firing can facilitate or depress transmission over tens to hundreds of milliseconds, timescales comparable to burst durations \cite{zucker2002short, jackman2017mechanisms}.
These dynamics reshape effective coupling without requiring large changes in baseline synaptic conductance.
Neuromodulators frequently target presynaptic release machinery, thereby tuning the timing and gain of transmission in ways that can strongly reshape circuit dynamics \cite{NadimBucher2014}.
In \textit{Tritonia diomedea}, serotonergic interneurons within the swim CPG show this principle directly: spike timing--dependent facilitation of their own synapses grants the network graded frequency control without an external modulatory system \cite{katz1994dynamic, sakurai2003spike}.

In \textit{Dendronotus iris}, experiments have characterized the swim CPG's connectivity, command neurons, and responses to targeted perturbations \cite{SakuraiKatz2016, sakurai2019command}.
The command-like Si1 neurons both initiate swimming and modulate swim frequency, providing the unified control signal that the circuit architecture must accommodate.

Here we build a computational model of the \textit{Dendronotus} swim CPG constrained by its known circuit topology, cellular properties, and synaptic kinetics.
Our analysis reveals three principal findings.
First, rhythm generation depends on contralateral E/I modules operating near an oscillatory threshold: ``almost-oscillatory'' subcircuits whose slow synaptic feedback creates the temporal framework for bursting.
Second, this distributed architecture confers resilience: the model reproduces experimental perturbation responses without parameter adjustment, confirming that no single neuron is essential for rhythm maintenance.
Third, neuromodulatory control of presynaptic release kinetics enables smooth frequency tuning, whereas postsynaptic conductance modulation produces abrupt, threshold-like transitions.
These results establish that network hysteresis, emergent from slow synaptic dynamics rather than cellular bistability, provides a robust yet flexible substrate for CPG function---reconnecting with insights from the early characterization of network oscillators while providing new mechanistic detail on how such circuits achieve controllability.
